\section{Classical principles and negative translation}\label{sec:negative}

\subsection{Exactly where the classical strength enters}

\begin{lemma}[Double negation and stable conclusions]\label{thm:stable-case}
\IPC\ proves $A\to\neg\neg A$ and $\neg\neg(A\vee\neg A)$.
Every formula $\neg\neg C$ is stable, meaning that
$\neg\neg\neg\neg C\to\neg\neg C$ is provable.
If $C$ is stable, the rule
\[
 \frac{\Gamma,A\vdash_I C\qquad\Gamma,\neg A\vdash_I C}
      {\Gamma\vdash_I C}
\]
is derivable with a fixed proof template using a supplied stability proof.
\end{lemma}
\begin{proof}
Under $A$ and $A\to\bot$, application gives $\bot$, proving the first
assertion. To prove the second, assume $h:\neg(A\vee\neg A)$. Assuming
$A$ and injecting it contradicts $h$, so $\neg A$ follows. Injecting that
negation again contradicts $h$, and discharging $h$ gives the result.
Stability of double negations is the triple-negation law from
Theorem~\ref{thm:int-laws} applied to $\neg C$.

For the rule, suppose $\neg C$. Under an additional $A$, the first branch
gives $C$, so $\neg A$ follows. The second branch now gives $C$, another
contradiction. Discharging $\neg C$ gives $\neg\neg C$; the supplied
stability proof gives $C$. No case split on the truth of $A$ occurred.
\end{proof}

\begin{theorem}[Equivalent classical extensions]\label{thm:classical-schemes}
Over \IPC, adding any one of the following \emph{schemes} gives exactly
\CPC:
\[
 A\vee\neg A,\qquad \neg\neg A\to A,\qquad
 ((A\to B)\to A)\to A.
\]
The equivalences are witnessed by primitive recursive proof translations.
\end{theorem}
\begin{proof}
Excluded middle eliminates double negations: its $A$ branch gives $A$,
while its $\neg A$ branch contradicts $\neg\neg A$ and gives $A$ by
falsehood elimination. Double-negation elimination gives excluded middle
by applying it to the second assertion of Lemma~\ref{thm:stable-case}.
For Peirce's law, suppose $h:(A\to B)\to A$. Under $\neg A$, any
assumed $A$ yields falsehood and hence $B$, giving $A\to B$. Applying
$h$ then contradicts $\neg A$. Thus $\neg\neg A$, and double-negation
elimination yields $A$.

Conversely specialize Peirce's scheme to $B=\bot$. Given $\neg\neg A$,
an assumption $\neg A$ yields falsehood and hence $A$. Therefore
$\neg A\to A$, and that instance of Peirce gives $A$. These are fixed
templates for replacing each added-axiom node of a proof. The distinction
between a scheme and an isolated instance matters: an assertion about
all substitution instances is what the theorem proves.
\end{proof}

\subsection{A constructive contextual Glivenko theorem}

\begin{theorem}[Glivenko with finite hypotheses]\label{thm:glivenko}
There are primitive recursive transformations proving
\[
       \Gamma\vdash_C A
       \quad\Longleftrightarrow\quad
       \Gamma\vdash_I\neg\neg A.
\]
Correctness is provable in \ipra\ as a statement about the two proof
checkers. In particular a classically valid propositional formula has an
intuitionistic proof of its double negation.
\end{theorem}
\begin{proof}
From a classical proof obtain its finite Boolean consequence certificate
by soundness. Reuse the leaf construction in the proof of
Theorem~\ref{thm:cpc-complete}: for every full literal context $L_a$ it
produces an \IPC\ proof of $A$ from $\Gamma,L_a$. That construction uses
only the signed evaluation lemma and falsehood elimination, even when
a premise in $\Gamma$ is false on the row. Append double-negation
introduction to obtain $\neg\neg A$ at every leaf.

Merge the leaves in the same binary pattern, but use the stable-conclusion
rule of Lemma~\ref{thm:stable-case}, with $C=\neg\neg A$, in place of
classical case analysis. It discharges the two opposite literals without
excluded middle. The final proof has precisely the context $\Gamma$ and
conclusion $\neg\neg A$. This is a bounded proof-table construction.
In the opposite direction, regard the intuitionistic proof as a classical
one and append the classical double-negation-elimination template.
\end{proof}

\begin{corollary}[Conservativity for stable conclusions]\label{thm:stable-conservativity}
If a supplied \IPC\ proof establishes $\neg\neg A\to A$, then any
\CPC\ proof of $A$ from $\Gamma$ can be changed into an \IPC\ proof of
$A$ from $\Gamma$. In particular this holds for negated conclusions.
\end{corollary}
\begin{proof}
Apply contextual Glivenko and then the stability proof. A negation is
stable by the triple-negation law.
\end{proof}

\subsection{A structural negative translation}

The truth-table method above gives an especially simple existence proof.
There is also a translation which follows the input proof structurally.
Define
\[
\begin{array}{ll}
 (p_i)^N=\neg\neg p_i,&\bot^N=\bot,\\
 (A\wedge B)^N=A^N\wedge B^N,&
 (A\vee B)^N=\neg\neg(A^N\vee B^N),\\
 (A\to B)^N=A^N\to B^N.&
\end{array}
\]
This is a primitive recursive formula transformation. Its output size is
bounded by a constant times the input size because each constructor adds
only a fixed number of nodes.

\begin{lemma}[Stability of translated formulas]\label{thm:negative-stable}
There is a primitive recursive constructor for \IPC\ proofs of
$\neg\neg A^N\to A^N$. If $Y$ is stable, then from $\neg\neg X$ and
$X\to Y$ one can derive $Y$ intuitionistically.
\end{lemma}
\begin{proof}
For the last assertion, an assumption $\neg Y$ and $X\to Y$ give
$\neg X$, contradicting $\neg\neg X$. This yields $\neg\neg Y$, hence $Y$.
For stability use structural induction. Double negations, including the
translated atom and disjunction cases, are stable. Falsehood is stable:
$\neg\neg\bot$ applied to the identity $\bot\to\bot$ gives $\bot$.
If $P,Q$ are stable, then $\neg\neg(P\wedge Q)$ gives $\neg\neg P$ and
$\neg\neg Q$ by composing hypothetical negations with projections.
Stability gives the components and hence their conjunction.
If $Q$ is stable, assume $\neg\neg(P\to Q)$ and $P$. Under $\neg Q$,
any function $P\to Q$ gives a contradiction, so $\neg(P\to Q)$ follows,
contradicting the first assumption. Hence $\neg\neg Q$ and then $Q$.
Discharge $P$. These are recursive proof templates for every constructor.
\end{proof}

\begin{theorem}[Negative translation of classical proofs]\label{thm:negative-translation}
A primitive recursive transformation maps every \CPC\ proof of $A$ from
$\Gamma$ to an \IPC\ proof of $A^N$ from $\Gamma^N$. Also \CPC\ proves
$A\leftrightarrow A^N$, uniformly by a primitive recursive proof constructor.
\end{theorem}
\begin{proof}
Translate the proof tree. Conjunction and implication rules commute with
the translation, as do assumptions and falsehood elimination. Disjunction
introduction first introduces $A^N\vee B^N$ and then its double negation.
For disjunction elimination the translated branches give
$(A^N\vee B^N)\to C^N$. Its principal premise is the double negation of
that disjunction; apply Lemma~\ref{thm:negative-stable}, since $C^N$ is stable.
The translation of an excluded-middle axiom is
$\neg\neg(A^N\vee\neg A^N)$, already proved in
Lemma~\ref{thm:stable-case}. Thus every translated rule has an intuitionistic
template. For the classical equivalence, evaluate both sides under a
Boolean valuation: double negation changes no bit, so a structural induction
identifies the values. Classical completeness supplies the equivalence
proof, or the classical double-negation templates give it directly.
\end{proof}

\subsection{Substitution into arithmetic without a reflection mistake}

\begin{proposition}[Decidable arithmetic instances]\label{thm:decidable-instances}
Suppose a finite collection of arithmetic formulas $P_i$ is supplied,
together with iPRA derivations of $P_i\vee\neg P_i$. A classical
propositional derivation whose atoms lie in this collection can be
translated primitive recursively into an iPRA derivation of its arithmetic
instance from the instantiated premises. Every bounded arithmetic formula
is eligible because of Theorem~\ref{thm:bounded-reflection}.
\end{proposition}
\begin{proof}
Decidability is preserved by every propositional connective. For
conjunction and disjunction inspect the two given decisions. For an
implication, a decision of its antecedent and consequent either provides
the implication, by a true consequent or a false antecedent, or provides
its negation, by a true antecedent and a false consequent. Falsehood is
decidable. Hence recursion constructs a decision proof for each formula
whose excluded-middle axiom occurs in the input derivation. Substitute
the arithmetic formulas throughout and replace those classical nodes by
these constructed proofs; every remaining propositional rule is already
an iPRA logical rule. This proves a syntactic proof-translation claim.
It does not supply a universal arithmetic truth predicate, nor does it
turn a proof-checker acceptance statement into unrestricted truth of an
arbitrary coded arithmetic formula.
\end{proof}

\section{What finite completeness does and does not assert}\label{sec:limits}

\subsection{Finite compactness statements that are proved here}

Every formula uses finitely many atoms, every proof uses finitely many
assumptions, and the proof constructors above use finite contexts. In
particular, a coded proof from an enumerated collection of hypotheses
comes with a finite support: traverse its hypothesis leaves and collect
their finitely many indices. This extraction is primitive recursive.
One can then rebuild the proof over precisely that finite subcollection.

There is also a useful finite-vocabulary observation. For $n$ fixed atoms
there are $2^{2^n}$ distinct Boolean truth tables, and every formula is
equivalent to one of them. Any \emph{finite} family of formulas over those
atoms may therefore be reduced to at most $2^{2^n}$ representative formulas
without changing its models, by evaluating and deduplicating its truth
tables. The representative indices and equivalence proofs are bounded
search outputs. This does not by itself provide a bound on when a new
truth table might appear in an infinite enumeration.

\subsection{Why arbitrary countable compactness is a different assertion}

An unrestricted claim that every finitely satisfiable countable theory has
an infinite satisfying assignment quantifies over an infinite object that
is not among the numerical finite certificates used in this paper. The
obstruction is concrete, not just terminological.

\begin{proposition}[Compactness would supply paths through binary trees]
\label{thm:compactness-tree}
Let $T$ be a prefix-closed binary tree with a node at every finite depth,
specified by a fixed primitive recursive membership predicate. From $T$
one can primitive recursively enumerate a countable propositional
theory every finite subcollection of which is satisfiable, such that any
infinite satisfying valuation determines an infinite path in $T$.
\end{proposition}
\begin{proof}
Use atoms $p_0,p_1,\ldots$. For each binary string
$s=(s_0,\ldots,s_{k-1})$ absent from $T$, include the clause
\[
 C_s=\bigvee_{i<k}\begin{cases}p_i&s_i=0,\\\neg p_i&s_i=1.\end{cases}
\]
It is false on exactly the valuations whose first $k$ bits are $s$.
For a total enumeration, scan all binary strings in length order and output
$\top$ for a string in $T$, or $C_s$ otherwise. This avoids any search for
a next missing string. For a finite subcollection let $L$ be the maximum
excluded-string length occurring there, with $L=0$ for the empty
subcollection. Choose a node $t\in T$ of length $L$ by bounded search among the
$2^L$ strings, using the assumed existence at that level. Its bits satisfy
every listed clause: otherwise an excluded string would be a prefix of
$t$, contrary to prefix closure. Thus finite satisfiability has explicit
witnesses at each finite stage. If a single infinite assignment satisfies
all clauses, none of its finite prefixes is excluded, so every prefix lies
in $T$. It is exactly an infinite path. For a tree supplied only with a
general computable decider, the same construction is effective relative
to that decider, but no primitive recursive running bound is asserted.
\end{proof}

Thus arbitrary countable compactness would solve an infinite path-selection
problem, usually expressed in second-order settings by a weak K\"onig
principle. We do not add such a principle to iPRA, do not claim a universal
primitive recursive path selector, and do not identify stage-by-stage
finite witnesses with a coherent infinite witness. This proof also pinpoints
why the finite maximal-extension construction in the intuitionistic section
is harmless: its entire universe has the explicit size $m$, and it performs
only $m$ greedy tests.

\subsection{Other boundaries of the result}

The Boolean and finite Kripke completeness theorems are full propositional
theorems for arbitrary finite input syntax, not bounds on a preselected
finite list of formulas. Nevertheless they should not be confused with
completeness for the ``intended meaning'' of all constructive arithmetic
propositions, with first-order completeness, or with arithmetic reflection.
For the intended-interpretation distinction see the original analysis by
Burgess \cite{burgess1981}; no additional choice-sequence principle from
that work is used here.

Our proof-system translations do not give polynomial simulations between
all proof systems. Truth tables and reconstructed cut-free proofs may be
exponentially large. The intuitionistic saturation table also deliberately
uses an exponential number of sequents. Nothing here settles $P$ versus
$NP$, gives a polynomial-time general SAT solver, or establishes a sharp
complexity bound for intuitionistic validity. Modal, quantified, and
first-order logics require separate syntax and metatheory and are not
silently included by the word ``logic''.
